% Options for packages loaded elsewhere
\PassOptionsToPackage{unicode}{hyperref}
\PassOptionsToPackage{hyphens}{url}
%
\documentclass[
]{article}
\usepackage{amsmath,amssymb}
\usepackage{lmodern}
\usepackage{iftex}
\ifPDFTeX
  \usepackage[T1]{fontenc}
  \usepackage[utf8]{inputenc}
  \usepackage{textcomp} % provide euro and other symbols
\else % if luatex or xetex
  \usepackage{unicode-math}
  \defaultfontfeatures{Scale=MatchLowercase}
  \defaultfontfeatures[\rmfamily]{Ligatures=TeX,Scale=1}
\fi
% Use upquote if available, for straight quotes in verbatim environments
\IfFileExists{upquote.sty}{\usepackage{upquote}}{}
\IfFileExists{microtype.sty}{% use microtype if available
  \usepackage[]{microtype}
  \UseMicrotypeSet[protrusion]{basicmath} % disable protrusion for tt fonts
}{}
\makeatletter
\@ifundefined{KOMAClassName}{% if non-KOMA class
  \IfFileExists{parskip.sty}{%
    \usepackage{parskip}
  }{% else
    \setlength{\parindent}{0pt}
    \setlength{\parskip}{6pt plus 2pt minus 1pt}}
}{% if KOMA class
  \KOMAoptions{parskip=half}}
\makeatother
\usepackage{xcolor}
\IfFileExists{xurl.sty}{\usepackage{xurl}}{} % add URL line breaks if available
\IfFileExists{bookmark.sty}{\usepackage{bookmark}}{\usepackage{hyperref}}
\hypersetup{
  pdftitle={Medición de la Interdisciplinariedad con Datos Internos Universitarios: Guía Operativa},
  pdfauthor={A. Rivero},
  hidelinks,
  pdfcreator={LaTeX via pandoc}}
\urlstyle{same} % disable monospaced font for URLs
\usepackage{longtable,booktabs,array}
\usepackage{calc} % for calculating minipage widths
% Correct order of tables after \paragraph or \subparagraph
\usepackage{etoolbox}
\makeatletter
\patchcmd\longtable{\par}{\if@noskipsec\mbox{}\fi\par}{}{}
\makeatother
% Allow footnotes in longtable head/foot
\IfFileExists{footnotehyper.sty}{\usepackage{footnotehyper}}{\usepackage{footnote}}
\makesavenoteenv{longtable}
\setlength{\emergencystretch}{3em} % prevent overfull lines
\providecommand{\tightlist}{%
  \setlength{\itemsep}{0pt}\setlength{\parskip}{0pt}}
\setcounter{secnumdepth}{-\maxdimen} % remove section numbering
\ifLuaTeX
  \usepackage{selnolig}  % disable illegal ligatures
\fi

\title{Medición de la Interdisciplinariedad con Datos Internos
Universitarios: Guía Operativa}
\author{A. Rivero}
\date{2026}

\begin{document}
\maketitle
\begin{abstract}
Las universidades poseen datos internos ricos --- repositorios de
publicaciones, registros de coautoría, historiales de codirección
doctoral y bases de datos de proyectos --- a los que los servicios
bibliométricos externos no tienen acceso. Presentamos un protocolo
operativo de medición que aprovecha estos datos institucionales para
evaluar investigación interdisciplinaria a nivel departamental e
institucional. El protocolo combina un panel bibliométrico de tres
componentes (diversidad, coherencia y efecto transcampo) con indicadores
solo institucionales (diversidad de coautoría, diversidad de
codirección, diversidad de paneles de financiación). Con un escenario
departamental simulado, mostramos cómo los datos internos permiten
distinguir amplitud polímata de integración genuina --- una diferencia
que las bases externas no capturan bien. Identificamos el principal
obstáculo de datos (citas para el efecto transcampo) y proponemos una
alternativa de dos componentes para instituciones sin suscripciones
bibliométricas comerciales.
\end{abstract}

\hypertarget{introducciuxf3n}{%
\section{Introducción}\label{introducciuxf3n}}

Las universidades buscan cada vez más fomentar y evaluar investigación
interdisciplinaria dentro de su planeación estratégica y sus sistemas de
aseguramiento de calidad. Sin embargo, las bases bibliométricas
comerciales (Web of Science, Scopus) ofrecen una visión parcial de la
actividad investigadora y, además, sus costos de suscripción pueden ser
altos. A esto se suma que esas bases no observan procesos
institucionales clave --- coautorías, codirecciones doctorales,
colaboraciones en proyectos y adscripciones departamentales --- que
muestran cómo ocurre realmente la integración del conocimiento.

Esta nota presenta una guía operativa para medir interdisciplinariedad
con predominio de datos internos. El enfoque central combina un panel
bibliométrico de tres componentes --- diversidad (\(\Delta\)),
coherencia (\(S\)) y efecto transcampo (\(E\)), desarrollado en la
revisión complementaria (Rivero, 2026) --- con indicadores
institucionales que capturan estructuras de colaboración. Mostramos el
protocolo en un escenario departamental simulado y añadimos
recomendaciones de extracción de datos, control de calidad e
interpretación.

\hypertarget{panorama-de-datos}{%
\section{Panorama de Datos}\label{panorama-de-datos}}

\hypertarget{quuxe9-tiene-una-universidad}{%
\subsection{Qué Tiene una
Universidad}\label{quuxe9-tiene-una-universidad}}

La mayoría de universidades intensivas en investigación mantienen:

\begin{itemize}
\tightlist
\item
  \textbf{Repositorios de publicaciones}: repositorios institucionales o
  CRIS (Current Research Information Systems) con metadatos de
  producción científica (autores, referencias, palabras clave).
\item
  \textbf{Sistemas de personal}: adscripción departamental de todo el
  personal académico.
\item
  \textbf{Registros doctorales}: información de dirección y codirección
  de tesis, incluyendo la afiliación de codirectores.
\item
  \textbf{Bases de proyectos}: registros de financiación con fuente,
  panel de evaluación y coinvestigadores.
\end{itemize}

Estas fuentes son completas dentro de la institución, pero no están
normalmente disponibles en servicios externos.

\hypertarget{quuxe9-suele-faltar}{%
\subsection{Qué Suele Faltar}\label{quuxe9-suele-faltar}}

\begin{itemize}
\tightlist
\item
  \textbf{Clasificación disciplinar de referencias}: asignación de
  referencias a categorías (por ejemplo, categorías Web of Science).
\item
  \textbf{Matriz de similitud entre categorías}: necesaria para computar
  diversidad Rao-Stirling. Puede derivarse de flujos de citación o
  estimarse con modelos de lenguaje (Cantone, 2025).
\item
  \textbf{Datos de citación}: quién cita cada publicación y desde qué
  campos, necesarios para \(E\).
\end{itemize}

La principal brecha operativa es el acceso a citación para calcular
\(E\). Por eso incluimos una alternativa de dos componentes.

\hypertarget{protocolo-de-mediciuxf3n}{%
\section{Protocolo de Medición}\label{protocolo-de-mediciuxf3n}}

\hypertarget{paso-1-panel-bibliomuxe9trico}{%
\subsection{Paso 1: Panel
Bibliométrico}\label{paso-1-panel-bibliomuxe9trico}}

Para cada investigador, calcular:

\textbf{Diversidad} (\(\Delta\)): índice Rao-Stirling sobre referencias:
\[\Delta = \sum_{i \neq j} d_{ij}\, p_i\, p_j, \qquad d_{ij} = 1 - s_{ij}\]
donde \(p_i\) es la proporción de referencias en la categoría \(i\),
\(s_{ij}\) es la similitud entre categorías y \(d_{ij}\) la distancia
correspondiente.

\textbf{Coherencia} (\(S\)): acoplamiento bibliográfico medio entre
publicaciones:
\[S = \frac{1}{\binom{n}{2}} \sum_{k < l} \cos(\mathbf{r}_k, \mathbf{r}_l)\]
donde \(\mathbf{r}_k\) es el vector de referencias de la publicación
\(k\).

\textbf{Efecto transcampo} (\(E\)): fracción de citas recibidas desde
fuera de la categoría principal del investigador.

\hypertarget{paso-2-indicadores-solo-institucionales}{%
\subsection{Paso 2: Indicadores Solo
Institucionales}\label{paso-2-indicadores-solo-institucionales}}

Además del panel bibliométrico, calcular:

\textbf{Diversidad de coautoría}:
\[\text{CoAuth} = \frac{\text{publicaciones con coautores de otros departamentos}}{\text{publicaciones totales}}\]

\textbf{Diversidad de codirección}:
\[\text{CoSup} = \frac{\text{codirecciones interdepartamentales}}{\text{direcciones totales}}\]

\textbf{Diversidad de paneles de financiación}: número de paneles
distintos en los que el investigador obtiene proyectos.

\hypertarget{paso-3-interpretaciuxf3n}{%
\subsection{Paso 3: Interpretación}\label{paso-3-interpretaciuxf3n}}

La combinación de indicadores bibliométricos e institucionales permite
separar integración de amplitud polímata:

\begin{longtable}[]{@{}
  >{\raggedright\arraybackslash}p{(\columnwidth - 10\tabcolsep) * \real{0.1667}}
  >{\raggedright\arraybackslash}p{(\columnwidth - 10\tabcolsep) * \real{0.1296}}
  >{\raggedright\arraybackslash}p{(\columnwidth - 10\tabcolsep) * \real{0.1296}}
  >{\raggedright\arraybackslash}p{(\columnwidth - 10\tabcolsep) * \real{0.1481}}
  >{\raggedright\arraybackslash}p{(\columnwidth - 10\tabcolsep) * \real{0.1296}}
  >{\raggedright\arraybackslash}p{(\columnwidth - 10\tabcolsep) * \real{0.2963}}@{}}
\toprule
\begin{minipage}[b]{\linewidth}\raggedright
Perfil
\end{minipage} & \begin{minipage}[b]{\linewidth}\raggedright
\(\Delta\)
\end{minipage} & \begin{minipage}[b]{\linewidth}\raggedright
S
\end{minipage} & \begin{minipage}[b]{\linewidth}\raggedright
CoAuth
\end{minipage} & \begin{minipage}[b]{\linewidth}\raggedright
CoSup
\end{minipage} & \begin{minipage}[b]{\linewidth}\raggedright
Interpretación
\end{minipage} \\
\midrule
\endhead
Integrador & Alta & Moderada-alta & Alta & Alta & Integración
transversal en publicaciones y procesos \\
Polímata & Alta & Baja & Baja & Baja & Amplitud de referencias sin
integración colaborativa \\
Especialista & Baja & Alta & Baja & Baja & Trayectoria focalizada
disciplinar \\
\bottomrule
\end{longtable}

\textbf{Idea clave:} una persona con \(\Delta\) alta y gran diversidad
de paneles de financiación puede parecer interdisciplinaria en datos
externos, pero si CoAuth y CoSup son nulos, el patrón es amplitud sin
integración.

Esto coincide con la distinción entrada/salida de la revisión principal:
la diversidad de insumos puede reflejar interdisciplinariedad
integrativa o solo yuxtaposición multidisciplinaria; para distinguirlas
se requieren coherencia e indicadores de proceso.

\hypertarget{demostraciuxf3n-departamento-simulado}{%
\subsection{Demostración: Departamento
Simulado}\label{demostraciuxf3n-departamento-simulado}}

Consideramos un departamento pequeño de Física y Ciencia de Materiales
con 3 investigadores. Usamos cinco categorías estilo Web of Science y
una matriz de similitud ilustrativa.

\begin{longtable}[]{@{}ll@{}}
\toprule
ID & Categoría \\
\midrule
\endhead
C1 & Física de la materia condensada \\
C2 & Ciencia de materiales \\
C3 & Química física \\
C4 & Óptica \\
C5 & Ingeniería eléctrica \\
\bottomrule
\end{longtable}

\begin{longtable}[]{@{}lrrrrr@{}}
\toprule
& C1 & C2 & C3 & C4 & C5 \\
\midrule
\endhead
C1 & 1.00 & 0.60 & 0.40 & 0.35 & 0.30 \\
C2 & 0.60 & 1.00 & 0.50 & 0.25 & 0.40 \\
C3 & 0.40 & 0.50 & 1.00 & 0.30 & 0.20 \\
C4 & 0.35 & 0.25 & 0.30 & 1.00 & 0.45 \\
C5 & 0.30 & 0.40 & 0.20 & 0.45 & 1.00 \\
\bottomrule
\end{longtable}

Los valores son ilustrativos; en práctica deben estimarse con datos
reales.

\textbf{Dra. Emma} (\(\mathbf{p}_E = (0.40, 0.30, 0.25, 0.03, 0.02)\)):
- Panel: \(\Delta = 0.42\), \(S = 0.55\), \(E = 0.22\) - Institucional:
CoAuth = 0.20, CoSup = 0.20, Paneles = 2 - \textbf{Perfil}: Integradora.

\textbf{Dr.~Farid} (\(\mathbf{p}_F = (0.25, 0.25, 0.25, 0.20, 0.05)\)):
- Panel: \(\Delta = 0.58\), \(S = 0.05\), \(E = 0.08\) - Institucional:
CoAuth = 0.00, CoSup = 0.00, Paneles = 4 - \textbf{Perfil}: Polímata no
integrativo.

\textbf{Dra. Greta} (\(\mathbf{p}_G = (0.70, 0.25, 0.03, 0.01, 0.01)\)):
- Panel: \(\Delta = 0.28\), \(S = 0.75\), \(E = 0.08\) - Institucional:
CoAuth = 0.00, CoSup = N/A, Paneles = 1 - \textbf{Perfil}: Especialista
en etapa temprana.

\hypertarget{guuxeda-de-implementaciuxf3n}{%
\section{Guía de Implementación}\label{guuxeda-de-implementaciuxf3n}}

\hypertarget{extracciuxf3n-de-datos}{%
\subsection{Extracción de Datos}\label{extracciuxf3n-de-datos}}

\begin{enumerate}
\def\labelenumi{\arabic{enumi}.}
\tightlist
\item
  \textbf{Publicaciones}: exportar del repositorio/CRIS con referencias
  completas.
\item
  \textbf{Afiliaciones de coautores}: integrar con sistema de RR. HH.
\item
  \textbf{Codirecciones doctorales}: integrar con registros de posgrado.
\item
  \textbf{Proyectos}: mapear proyectos a taxonomías de paneles.
\item
  \textbf{Citación}: si existe, importar desde OpenAlex, WoS o Scopus
  para \(E\).
\end{enumerate}

\hypertarget{controles-de-calidad}{%
\subsection{Controles de Calidad}\label{controles-de-calidad}}

\begin{itemize}
\tightlist
\item
  \textbf{Cobertura de categorías en referencias}: si falta
  \textgreater25\%, \(\Delta\) será inestable (Nakhoda et al., 2023).
\item
  \textbf{Carreras tempranas}: para \textless5 años pos-PhD o
  \textless10 publicaciones, reportar intervalos de confianza junto con
  estimaciones puntuales.
\item
  \textbf{Calibración y ambigüedad}: umbrales ilustrativos deben
  calibrarse con distribuciones locales; si un intervalo cruza fronteras
  de perfil, clasificar como ``ambiguo'' y solicitar revisión
  cualitativa.
\item
  \textbf{Cambios de afiliación}: estratificar por periodo para evitar
  inflar CoAuth artificialmente.
\end{itemize}

\hypertarget{reporte}{%
\subsection{Reporte}\label{reporte}}

Reportar en tres niveles:

\begin{enumerate}
\def\labelenumi{\arabic{enumi}.}
\tightlist
\item
  \textbf{Individual}: perfil por investigador (\(\Delta\), S, E,
  CoAuth, CoSup, Paneles).
\item
  \textbf{Departamental}: distribuciones y medianas.
\item
  \textbf{Institucional}: agregados comparativos entre áreas.
\end{enumerate}

\textbf{No usar ranking compuesto único.} El panel es multidimensional y
debe interpretarse como vector de evidencia (Rafols, 2019).

\hypertarget{alternativa-sin-datos-de-citaciuxf3n}{%
\subsection{Alternativa sin Datos de
Citación}\label{alternativa-sin-datos-de-citaciuxf3n}}

Si no hay datos de citación, usar panel interno de dos componentes
(\(\Delta\), S) + indicadores institucionales (CoAuth, CoSup, Paneles):

\begin{longtable}[]{@{}llllll@{}}
\toprule
Perfil & \(\Delta\) & S & CoAuth & CoSup & Interpretación \\
\midrule
\endhead
Integrador & Alta & Moderada-alta & Alta & Alta & Integración en insumos
y procesos \\
Polímata & Alta & Baja & Baja & Baja & Amplitud desconectada \\
Especialista & Baja & Alta & Baja & Baja & Foco disciplinar \\
\bottomrule
\end{longtable}

Se pierde capacidad para medir impacto transcampo (\(E\)), pero se
mantiene capacidad para identificar integración.

\hypertarget{limitaciones-y-extensiones}{%
\section{Limitaciones y Extensiones}\label{limitaciones-y-extensiones}}

Primero, el protocolo asume datos limpios y estructurados; muchas
instituciones aún no los tienen. Segundo, caracteriza tipo de
interdisciplinariedad, no calidad intrínseca de resultados. Tercero, los
umbrales son orientativos y deben calibrarse empíricamente. Cuarto, en
portafolios pequeños la incertidumbre puede ser grande; por tanto, las
reglas de clasificación deben ser sensibles a intervalos y no solo a
valores puntuales.

Entre las extensiones posibles: indicadores semánticos (título/resumen),
indicadores docentes (docencia interdepartamental) y análisis temporal
de trayectorias.

\hypertarget{conclusiones}{%
\section{Conclusiones}\label{conclusiones}}

Las universidades pueden medir interdisciplinariedad con mayor
resolución si combinan panel bibliométrico y datos internos de
colaboración. El protocolo propuesto distingue integración genuina de
amplitud polímata, evita la reducción a un único número y permite
decisiones más auditables. El principal cuello de botella sigue siendo
la citación para \(E\), pero la alternativa interna (\(\Delta\), S +
indicadores institucionales) ofrece una base sólida para uso operativo.

\hypertarget{referencias}{%
\section{Referencias}\label{referencias}}

\begin{itemize}
\tightlist
\item
  Cantone, G. G. (2025). Estimation of disciplinary similarity with
  large language models. \emph{Scientometrics}, 130(10):5345--5373.
\item
  Nakhoda, M., Whigham, P., and Zwanenburg, S. (2023). Quantifying and
  addressing uncertainty in the measurement of interdisciplinarity.
  \emph{Scientometrics}, 128:6107--6127.
\item
  Rafols, I. (2019). S\&T indicators in the wild: contextualization and
  participation for responsible metrics. \emph{Research Evaluation},
  28(1):7--22.
\item
  Rivero, A. (2026). Measuring interdisciplinarity: A multi-component
  indicator panel for research evaluation. Manuscript.
\end{itemize}

\end{document}
